%% Steven F. Wolf, Ph.D. - Comprehensive CV
%% Use as master reference when tailoring role-specific CVs.
%%
%% Compile with: pdflatex main_example.tex

\documentclass[10pt,letterpaper,ragged2e,withhyper,normalphoto]{altacv}
\usepackage{sfwCV}
\usepackage{multicol}

\begin{document}
\name{Steven F.\ Wolf Ph.D.}

\tagline{Data Scientist | Research Scientist | Technical Leader}

\makecvheader

\cvsection{Professional Summary}

Data Scientist with a Ph.D. in Physics and 10+ years of experience in
statistical modeling, computational analysis, and cross-functional
research leadership. Proven track record in securing \$1.3M in federal
grants and delivering tools used across departments. Combines deep
domain expertise in numerical methods and approximation techniques
with strong Python/R development skills and experience translating
complex problems for diverse stakeholders.

\cvsection{Technical Expertise}

\begin{description}
\item[Programming \& ML:] Python (NumPy, Pandas, Matplotlib,
  TensorFlow, scikit-learn), R (statistical modeling, Shiny, Plotly),
  SQL, Git
\item[Statistical \& Computational:] Statistical Models, Numerical
  Methods, Approximation Techniques, Complex Problem Solving, FEM/FEA
\item[Data \& Infrastructure:] ETL Pipelines, Cluster Computing, Data
  Visualization, Dashboard Creation, Airflow
\item[Research \& Communication:] Cross-Functional Collaboration,
  Technical Reporting, Program Management, Mentoring, Stakeholder
  Communication
\item[Tools \& Platforms:] Git, GitHub, LaTeX, Bash Scripting,
  HTML/CSS, JavaScript
\end{description}

\cvsection{Professional Experience}

\cvevent{Assistant Professor / Research Group Lead}{East Carolina
  University}{August 2015 -- May 2025}{Greenville, NC}
\begin{description}
\item[Financial Strategy:] Secured >\$1.3M in federal and
  institutional grants as PI (NSF-IUSE: \$900K; other sources)
\item[Program Leadership:] Led ECU Learning Assistant Program through
  planning, implementation, and evaluation; secured permanent funding
  from Dean's office (\$500K+)
\item[Curriculum Development:] Designed and taught introductory
  physics labs: developed and validated assessment of scientific
  practices; coordinated assessment of scientific practices across
  multiple STEM disciplines; mentored graduate TAs; curriculum adopted
  by Physics Department; still in use post-departure
\item[Research Output:] Published 25 peer-reviewed research papers;
  served as PI on multiple federal and state grants
\end{description}

\divider

\cvevent{Physics Lecturer / Mentor}{Texas State University}{August
  2014 -- August 2015}{San Marcos, TX}
\begin{itemize}
\item Taught college physics courses to classes of 30--100+ students
  using active, interactive learning techniques
\item Mentored 10 Learning Assistants and served as 1-1 mentor for the
  H-LSAMP program
\end{itemize}

\divider

\cvevent{Research Associate (Ph.D.)}{Michigan State University}{August
  2012 -- August 2014}{East Lansing, MI}
\begin{description}
\item[Data Analysis:] Created analysis and visualization tools with
  the R programming language, collaborating with research teams to
  address domain-specific questions (3 peer-reviewed publications)
\item[Project Management:] Conducted research independently, driving
  projects through all phases, including design, data collection, and
  delivery
\end{description}

\cvsection{Education}

\cvevent{Ph.D.\ in Physics}{Michigan State University}{2006 --
  2012}{East Lansing, MI} Thesis: Matter-wave amplification and
femtosecond laser spectroscopy. Additional research in expert-novice
cognitive structures.

\cvevent{M.S.\ in Physics}{Dartmouth College}{2003 -- 2005}{Hanover,
  NH} Research in solar wind phenomena and non-exponential decay.

\cvevent{B.S.\ in Physics and Mathematics}{Valparaiso University}{1999
  -- 2003}{Valparaiso, IN}

\cvsection{Projects}

\cvevent{AI Model Evaluation / Physics Domain Expert}{Handshake AI /
  Mercor Intelligence}{September 2025 -- Present}{Remote}
\begin{itemize}
\item Provide domain expertise in Physics and Data Science to improve
  model performance for top AI labs
\end{itemize}

\cvevent{XLABs -- Introductory Physics Lab Assessments} {East Carolina
  University}{2017 -- 2021}{Greenville, NC}

\begin{itemize}
\item Wrote laboratory curriculum for the introductory physics labs
  across the entire department, utilizing an Argument Driven Inquiry
  (ADI) framework to emphasize scientific practices. Curriculum is
  still in use by all instructors;
\item Designed and validated lab assessments targeting scientific
  practices across the introductory physics curriculum;
\item Trained Graduate TAs to implement the ADI-based curriculum in
  pedagogically appropriate ways, created materials that are used by
  current faculty to continue the pedagogical training;
\item Designed and carried out research investigating how students
  utilize scientific practices in the context of this lab curriculum
  as well as within the larger institutional transformation;
\end{itemize}

\cvevent{concorR R Package}{github.com/atraxlab/concorr}{September
  2019 -- August 2022}{} R package implementing CONCOR (CONvergence of
iterated CORrelations) for positional analysis of social networks,
grouping nodes by structural equivalence

\cvevent{PERC Proceedings Editorial Team}{Physics Education Research
  Conference}{May 2018 -- December 2020}{Remote}
\begin{description}
\item[Editor in Chief (Jan 2020 -- Dec 2020):] Led editorial team,
  coordinated delivery of conference proceedings during COVID
  pandemic; created fully virtual conference
\item[Editor (Jan 2019 -- Dec 2019):] Led peer review assignments;
  created timeline for double-confidential review process
\item[Associate Editor (May 2018 -- Dec 2018):] Coordinated
  proceedings delivery and maintained editorial standards
\end{description}

\cvsection{Selected Publications}

\begin{description}
\item[Wolf, S.F., Sault, T.M., Suda, T., \& Traxler, A.L.]
  (2022). Social Network Development in Classrooms. \textit{Applied
    Social Network Analysis} 7, 24.
\item[Wolf, S.F. \& Sprague, M.W.] (2022). Introductory Physics Labs:
  A Tale of Two Transformations. \textit{The Physics Teacher} 60, 372.
\item[Wolf, S.F., Dougherty, D.P., \& Kortemeyer, G.] (2012). Rigging
  the Deck. \textit{Phys. Rev. ST-PER} 8, 020116.
\item[Wolf, S.F., Dougherty, D.P., \& Kortemeyer, G.] (2012). An
  Empirical Approach to Interpreting Card-sorting
  Data. \textit{Phys. Rev. ST-PER} 8, 010124.
\end{description}

\cvsection{Certifications}

\begin{description}
\item[Supervised Machine Learning:] Regression and Classification
  (Coursera, Sep 2022)
\item[Advanced Learning Algorithms:] (Coursera, Oct 2022)
\item[Unsupervised Learning:] Recommenders, Reinforcement Learning
  (Coursera, Oct 2022)
\item[Linux Commands \& Shell Scripting:] (IBM, Dec 2022)
\item[Databases and SQL for Data Science:] with Python, with Honors
  (IBM, Jan 2023)
\item[ETL and Data Pipelines:] Shell, Airflow and Kafka (Coursera, Jul
  2023)
\end{description}

\cvsection{Awards}

\begin{description}
\item[AAPT Outstanding Teaching Assistant Award] -- American
  Association of Physics Teachers (2012)
\item[Lyman Briggs College Outstanding TA Award] (\$1000) -- Michigan
  State University (2012)
\item[CoGS Travel Grant Award] (\$500) -- Michigan State University
  (2011)
\item[Sigma Pi Sigma Induction] -- National Physics Honors Society
  (2003)
\item[Presidential Academic Honors for Student Athletes] -- Valparaiso
  University (2000--2003)
\end{description}


\end{document}
